\chapter{Quantization, ADCs, and DACs}

\chapterline

In the last chapter we cut time into a grid. Sampling answered the question \emph{when} do we
look at a signal, and Nyquist told us how often we must look so that nothing is lost. But there
is a second grid hiding inside every digital system, and it is just as important and far less
discussed: the grid on the \emph{vertical} axis. A computer cannot store the real number
$0.7182818\ldots$ It stores a finite pattern of bits. So when the analog-to-digital converter
(the ADC) measures a voltage, it must not only pick a moment in time -- it must also snap the
measured value to the nearest rung on a ladder of allowed values. That snapping is
\textbf{quantization}, and it is the subject of this chapter.

Sampling is about the horizontal grid; quantization is about the vertical grid. Together they
turn a smooth analog wiggle into a finite list of integers -- the only thing a microcontroller,
a CPU, or a DSP chip can actually hold. Going the other way, a digital-to-analog converter (the
DAC) reads integers back out and rebuilds a voltage you can send to a speaker or a motor. This
chapter builds the geometric picture of that vertical grid, derives the famous ``about 6~dB per
bit'' rule from scratch, opens up the two ADC architectures you will actually meet
(successive-approximation and sigma-delta), and shows why oversampling and \emph{adding noise on
purpose} (dithering) can buy you resolution that the raw bit count seems to forbid. We end at the
bench, reading a real sensor through the ADC on an Arduino and a Raspberry Pi.

\section{The vertical grid: from a value to a code}
\label{ch19:vertical-grid}

Picture the ADC's job geometrically. The input is a voltage $x$ somewhere in an allowed range,
say from $0$ to a reference voltage $V_{\mathrm{ref}}$. The output is one of a finite set of
integer \emph{codes}: $0, 1, 2, \ldots, 2^N - 1$ for an $N$-bit converter. The converter draws a
ladder of evenly spaced horizontal lines across the input range and reports which rung your
voltage is nearest. The spacing between rungs is the fundamental parameter governing
resolution and quantization noise:
\[
  \Delta \;=\; \frac{V_{\mathrm{ref}}}{2^N}.
\]
We call $\Delta$ the \textbf{quantization step} or the value of one \textbf{LSB} (least
significant bit) -- the smallest change in voltage that can possibly change the output code.

\begin{intuitionbox}
Sampling chops a continuous curve into a row of vertical pins (one per sample instant).
Quantization then files the top of each pin down to the nearest line on a piece of ruled paper.
After both operations the signal lives entirely on a finite lattice: finitely many time slots
across, finitely many value-rungs up. A digital signal is a set of points on graph paper, and
nothing in between exists.
\end{intuitionbox}

Figure~\ref{ch19:fig-staircase} shows the transfer characteristic -- output code versus input
voltage -- for a tiny 3-bit converter (eight rungs). It is a staircase. Inside each tread the
output does not change at all; cross a riser and the code jumps by one. The flat treads are
exactly where information is being thrown away: every input voltage in a tread is reported by the
same code.

\begin{figure}[H]
\centering
\begin{tikzpicture}
\begin{axis}[
  width=11cm, height=6cm,
  title={Quantizer transfer characteristic (3-bit, 8 codes)},
  title style={font=\small},
  xlabel={input voltage $x$ (in units of $V_{\mathrm{ref}}$)},
  ylabel={output code},
  xmin=0, xmax=1, ymin=0, ymax=8,
  axis lines=left,
  ytick={0,1,2,3,4,5,6,7},
  xtick={0,0.25,0.5,0.75,1},
  enlargelimits=false,
]
% ideal line
\addplot[gray, dashed, domain=0:1, samples=2] {8*x};
% staircase: codes 0..7, top code saturates at 7
\addplot[sigblue, thick, const plot, no marks] coordinates {
 (0,0) (0.125,1) (0.25,2) (0.375,3) (0.5,4) (0.625,5) (0.75,6) (0.875,7) (1,7)
};
\end{axis}
\end{tikzpicture}
\caption{The quantizer transfer characteristic for a 3-bit converter (8 codes). The dashed line
is the ideal continuous response; the blue staircase is what you actually get. Each flat tread
has width $\Delta = V_{\mathrm{ref}}/8$, and any input within a tread maps to the same code.}
\label{ch19:fig-staircase}
\end{figure}

\begin{definition}[Uniform quantizer]
\label{ch19:def-quantizer}
A \textbf{uniform quantizer} with step $\Delta$ maps a real input $x$ to the nearest multiple of
$\Delta$:
\[
  Q(x) \;=\; \Delta \cdot \mathrm{round}\!\left(\frac{x}{\Delta}\right).
\]
The integer $\mathrm{round}(x/\Delta)$, clipped to the available code range, is the digital code
stored in memory. Note the half-step alignment: $\mathrm{round}(x/\Delta)$ places its decision thresholds halfway between codes at $(k+\tfrac{1}{2})\Delta$, centering the error on $[-\Delta/2,+\Delta/2]$, whereas an unshifted hardware staircase like Figure~\ref{ch19:fig-staircase} switches on integer multiples $k\Delta$ ($\lfloor x/\Delta\rfloor$, truncation) unless shifted by $\Delta/2$.
\end{definition}

\paragraph{Bit depth.} The number $N$ is the \textbf{bit depth} or resolution. Each extra bit
doubles the number of rungs and halves $\Delta$. The vocabulary you will hear:
\begin{itemize}
  \item \textbf{8-bit}: 256 levels. Old samplers, simple sensors, control loops where you only
        need coarse values.
  \item \textbf{10/12-bit}: the ADCs built into microcontrollers (Arduino Uno is 10-bit,
        many Cortex-M parts are 12-bit). Plenty for an IMU axis or a slider.
  \item \textbf{16-bit}: CD-quality audio. 65{,}536 levels.
  \item \textbf{24-bit}: studio audio and high-end sigma-delta sensor front-ends.
\end{itemize}

\begin{trapbox}
Bit depth and sample rate are independent axes and people constantly confuse them. Sample rate
(Chapter~18) is the horizontal grid spacing and is governed by Nyquist; it decides the highest
\emph{frequency} you can represent. Bit depth is the vertical grid spacing; it decides the
smallest \emph{amplitude change} you can represent, which sets the noise floor. A 1~MHz, 4-bit
converter and a 1~kHz, 24-bit converter are good at completely different things. Always ask both
``how fast?'' and ``how fine?''
\end{trapbox}

\subsection*{Mid-rise versus mid-tread, and clipping}
There are two ways to position the staircase. A \emph{mid-tread} quantizer has a flat tread
centered on zero, so small inputs map to code zero -- handy because true silence reads as exactly
zero. A \emph{mid-rise} quantizer has a riser at zero, so it never outputs an exact zero but
treats positive and negative symmetrically. Audio converters are usually mid-tread. The
difference is one half-LSB of offset and rarely matters in practice, but it explains why your DC
reading might sit on $0$ or might dither between $-1$ and $+1$ codes.

What always matters is \textbf{clipping}. The ladder has a top and a bottom. Push the input above
$V_{\mathrm{ref}}$ (or below ground) and the code saturates at its maximum (or minimum) and
\emph{stays there}. Geometrically the staircase goes flat forever. Clipping acts as a hard wall
rather than a gentle rounding: it flattens the peaks of your waveform into plateaus, generating a
spray of harmonics. Sampling errors can be undone by a filter; clipping destroys information
irreversibly. Set your input gain so peaks land just below full scale, never above.

\begin{connectionbox}
The round-to-nearest rule is exactly the projection idea from the linear-algebra companion,
applied to a one-dimensional lattice. $Q(x)$ is the closest point of the set $\Delta\Z$ to $x$,
i.e.\ the orthogonal projection of $x$ onto a discrete grid. Quantization error is the residual
$x - Q(x)$, the part of $x$ that the grid cannot represent -- the same ``what is left over after
projection'' object you met with Fourier coefficients, now in the crudest possible space.
\end{connectionbox}

\section{Quantization error as noise}
\label{ch19:error-as-noise}

The interesting quantity is not $Q(x)$ but the \textbf{error} it makes:
\[
  e \;=\; Q(x) - x.
\]
Because $Q$ rounds to the nearest rung, and the rungs are spaced $\Delta$ apart, the error can
never exceed half a step in magnitude:
\[
  -\tfrac{\Delta}{2} \;\le\; e \;\le\; \tfrac{\Delta}{2}.
\]
That bound is exact and unconditional. The error is the vertical distance from each pin-top down
(or up) to the nearest ruled line, and that distance is at most half the line spacing.

One modeling assumption makes quantization tractable. For a busy signal -- music, an accelerometer
trace, anything that swings across many codes from sample to sample -- the error $e$ behaves like
a \emph{random number, uniformly spread} across the interval $[-\Delta/2, \Delta/2]$, and roughly
independent from one sample to the next. We do not need this to be exactly true; we need it to be
a good model, and for active signals it is excellent. Under that model, quantization is just a
weak, broadband \textbf{noise} added to the clean signal.

\begin{whybox}
Why is treating a \emph{deterministic} rounding rule as random legitimate? Because the signal
sweeps across the treads in a way that is, from the error's point of view, effectively
unpredictable. If $x$ moves by a fraction of an LSB that has no simple relation to $\Delta$, the
fractional part $x/\Delta \bmod 1$ tumbles through $[0,1)$ and visits all positions about
equally. The error is a deterministic function of an effectively-uniform phase, so it inherits a
uniform distribution. This is the same reason that a fast-spinning analog phase looks like a
random angle. The model fails only when the signal is too quiet or too slow to stir the
phase -- which is precisely the case dithering will fix in Section~\ref{ch19:dither}.
\end{whybox}

\subsection*{The power of the noise}
Model $e$ as uniform on $[-\Delta/2, \Delta/2]$. Its mean is zero, and its power is its variance
-- a textbook integral from the probability companion:
\[
  \sigma_e^2 \;=\; \E[e^2] \;=\; \frac{1}{\Delta}\int_{-\Delta/2}^{\Delta/2} e^2 \,\dd e
  \;=\; \frac{\Delta^2}{12}.
\]
The result $\Delta^2/12$ is used throughout DSP. It says the
quantization noise power depends only on the step size, and shrinks as the \emph{square} of the
step. Halve $\Delta$ (add one bit) and you cut the noise power by a factor of four.

\begin{connectionbox}
$\sigma_e^2 = \Delta^2/12$ is just the variance of a uniform distribution, derived in the
probability book. The factor of $1/12$ comes straight from integrating $u^2$: $\Var(U)$ for $U$ uniform on an interval of width $w$
is $w^2/12$, and here $w=\Delta$. Quantization is your first concrete meeting with a
\emph{random signal} -- a noise process whose statistics you can compute exactly. Chapter~20
picks up this thread and treats noise in earnest.
\end{connectionbox}

\begin{happeningbox}
Visualize what $\Delta^2/12$ means on the graph. After quantizing, your reconstructed waveform is
the true waveform plus a thin, grassy fuzz of height about $\pm\Delta/2$. That fuzz is spread
fairly evenly across \emph{all} frequencies (it is broadband, ``white''). It is not a distortion of
the shape so much as a fine hiss riding on top. Its total power is fixed by $\Delta$
alone -- it does not care how loud or quiet your signal is. This fact drives the
SNR analysis below: a loud signal stands far above a fixed hiss; a quiet one drowns in it.
\end{happeningbox}

\section{Signal-to-noise ratio and the 6~dB-per-bit rule}
\label{ch19:snr}

Now we can answer the question every spec sheet addresses: how good is an $N$-bit converter? We
compare the power of the largest signal it can carry against the power of the quantization hiss.
That ratio, in decibels, is the \textbf{signal-to-noise ratio} (SNR).

Take a full-scale sinusoid -- the loudest signal that fits without clipping. Its peak-to-peak
swing spans the full converter range $2^N \Delta$, so its amplitude is $A = 2^N\Delta/2 = 2^{N-1}\Delta$.
A sinusoid of amplitude $A$ has power
\[
  P_{\mathrm{sig}} \;=\; \frac{A^2}{2} \;=\; \frac{(2^{N-1}\Delta)^2}{2}
  \;=\; \frac{2^{2N}\,\Delta^2}{8}.
\]
The noise power is $\sigma_e^2 = \Delta^2/12$. Their ratio is
\[
  \mathrm{SNR} \;=\; \frac{P_{\mathrm{sig}}}{\sigma_e^2}
  \;=\; \frac{2^{2N}\Delta^2/8}{\Delta^2/12} \;=\; \frac{12}{8}\,2^{2N}
  \;=\; \frac{3}{2}\,2^{2N}.
\]
The $\Delta^2$ cancels, so the SNR of an ideal converter does not
depend on the reference voltage or the step size; it depends only on the number of bits. Now
convert to decibels:
\[
  \mathrm{SNR}_{\mathrm{dB}} \;=\; 10\log_{10}\!\left(\frac{3}{2}\,2^{2N}\right)
  \;=\; 10\log_{10}\tfrac{3}{2} \;+\; 2N\cdot 10\log_{10} 2.
\]
Plug in $10\log_{10}2 \approx 3.0103$ and $10\log_{10}\tfrac32 \approx 1.76$:
\[
  \boxed{\;\mathrm{SNR}_{\mathrm{dB}} \;\approx\; 6.02\,N \;+\; 1.76 \ \text{dB}.\;}
\]

\begin{focusbox}
\textbf{The 6~dB-per-bit rule.} Each additional bit of resolution buys you about $6.02$~dB of
SNR -- roughly a factor of two in amplitude, four in power. The $+1.76$~dB is the bonus from
using a full-scale sinusoid (a constant that depends on the signal's crest factor; it changes if
the signal is not a sinusoid, but the $6.02\,N$ slope is universal).
\end{focusbox}

\begin{example}[Putting numbers on it]
\label{ch19:ex-snr}
\begin{itemize}
  \item \textbf{8-bit:} $6.02(8) + 1.76 \approx 49.9$~dB. Audibly hissy for music, fine for a
        control signal.
  \item \textbf{12-bit} (typical MCU ADC): $6.02(12)+1.76 \approx 74$~dB. Excellent for an IMU
        axis or a slider.
  \item \textbf{16-bit} (CD): $6.02(16)+1.76 \approx 98$~dB. The hiss sits about 98~dB below a
        full-scale tone -- below the threshold of hearing in a normal room.
  \item \textbf{24-bit:} $\approx 146$~dB in theory. In practice the analog electronics
        (thermal noise in resistors, the reference, the op-amps) hit a wall around 20--21 bits
        of \emph{effective} resolution; the last bits are buried in real noise. This is why we
        speak of ENOB.
\end{itemize}
\end{example}

\begin{definition}[ENOB]
\label{ch19:def-enob}
The \textbf{effective number of bits} of a real converter is obtained by measuring its actual SNR
(including all real-world noise and distortion) and inverting the formula:
\[
  \mathrm{ENOB} \;=\; \frac{\mathrm{SNR}_{\mathrm{dB}} - 1.76}{6.02}.
\]
A converter labeled ``16-bit'' with a measured SNR of 90~dB has $\mathrm{ENOB}\approx 14.7$ bits.
The marketing number is the size of the register; ENOB is what you actually get.
\end{definition}

Figure~\ref{ch19:fig-snr} plots the ideal line. It is a straight line with slope $6.02$~dB/bit.

\begin{figure}[H]
\centering
\begin{tikzpicture}
\begin{axis}[
  width=11cm, height=5.5cm,
  title={Theoretical SNR vs.\ bit depth: $6.02N + 1.76$\,dB},
  title style={font=\small},
  xlabel={bit depth $N$},
  ylabel={ideal SNR (dB)},
  xmin=0, xmax=24, ymin=0, ymax=150,
  axis lines=left,
  xtick={0,4,8,12,16,20,24},
  ytick={0,25,50,75,100,125,150},
  enlargelimits=false,
  grid=both, grid style={gray!20},
]
\addplot[specred, thick, domain=0:24, samples=2] {6.02*x + 1.76};
\addplot[only marks, sigblue, mark=*] coordinates {(8,49.9)(12,74.0)(16,98.1)(24,146.2)};
\end{axis}
\end{tikzpicture}
\caption{Ideal SNR versus bit depth, $\mathrm{SNR}\approx 6.02N+1.76$~dB. Dots mark the common
depths from Example~\ref{ch19:ex-snr}. Real converters fall below this line by their noise and
distortion; the vertical gap is what ENOB captures.}
\label{ch19:fig-snr}
\end{figure}

\begin{walkbox}
If quantization noise power is fixed at $\Delta^2/12$ regardless of the signal, what happens to
the SNR of a \emph{quiet} signal -- say one using only the bottom few codes? Its $P_{\mathrm{sig}}$
plummets while the noise stays put, so SNR collapses. This is why recording engineers chase a hot
input level, and why a 16-bit recording made at $-40$~dB has the effective fidelity of a
roughly 9-bit recording. Use the full available dynamic range of the ladder.
\end{walkbox}

\section{Seeing and hearing quantization}
\label{ch19:seeing}

Let us make the hiss visible. Figure~\ref{ch19:fig-quantsine} quantizes a sine wave to a coarse
3 bits so the staircase is obvious. The smooth curve is the original; the stepped curve is what
you store; the gap between them, sample by sample, is the error $e$ -- and you can see it is a
little sawtooth-like wiggle bounded by $\pm\Delta/2$.

\begin{figure}[H]
\centering
\begin{tikzpicture}
\begin{axis}[
  width=12cm, height=5.5cm,
  title={Quantization of a continuous sine wave (3 bits)},
  title style={font=\small},
  xlabel={time $t$}, ylabel={amplitude},
  xmin=0, xmax=1, ymin=-1.25, ymax=1.25,
  axis lines=left, domain=0:1, samples=200,
  enlargelimits=false, legend pos=north east, legend cell align=left,
]
\addplot[sigblue, thick] {sin(deg(2*pi*2*x))};
\addlegendentry{original}
\addplot[specred, thick, samples=200] {0.25*round(sin(deg(2*pi*2*x))/0.25)};
\addlegendentry{3-bit quantized}
\end{axis}
\end{tikzpicture}
\caption{A sine wave quantized to 3 bits (step $\Delta=0.25$ on a $[-1,1]$ range). The red
staircase is the stored signal; the vertical gaps to the blue curve are the quantization error,
each no larger than $\Delta/2=0.125$.}
\label{ch19:fig-quantsine}
\end{figure}

The error itself is worth a plot of its own. Figure~\ref{ch19:fig-quanterror} draws
$e = Q(x) - x$ for that same 3-bit sine. It is a jagged little signal that lives entirely
inside the band $[-\Delta/2, \Delta/2]$ and never escapes it -- the bound from
Section~\ref{ch19:error-as-noise} made visible. For a busy signal this trace looks like grass;
that is the white-noise model at work.

\begin{figure}[H]
\centering
\begin{tikzpicture}
\begin{axis}[
  width=12cm, height=4.5cm,
  title={Quantization error waveform: $e(t) = Q(x(t)) - x(t)$},
  title style={font=\small},
  xlabel={time $t$}, ylabel={error $e$},
  xmin=0, xmax=1, ymin=-0.2, ymax=0.2,
  axis lines=left, domain=0:1, samples=200,
  enlargelimits=false,
  ytick={-0.125,0,0.125}, yticklabels={$-\Delta/2$,$0$,$\Delta/2$},
]
\addplot[gray, dashed, domain=0:1, samples=2] {0.125};
\addplot[gray, dashed, domain=0:1, samples=2] {-0.125};
\addplot[specred, thick, samples=200] {0.25*round(sin(deg(2*pi*2*x))/0.25) - sin(deg(2*pi*2*x))};
\end{axis}
\end{tikzpicture}
\caption{The quantization error $e = Q(x) - x$ for the 3-bit sine of
Figure~\ref{ch19:fig-quantsine}. It stays strictly between the dashed bounds $\pm\Delta/2$ and,
for an active signal, scatters like broadband noise.}
\label{ch19:fig-quanterror}
\end{figure}

\begin{labbox}
Quantize a tone yourself and \emph{listen} to the error. In Python:
\begin{lstlisting}[language=Python]
import numpy as np
fs = 44100
t  = np.arange(fs) / fs           # 1 second
x  = 0.9 * np.sin(2*np.pi*440*t)  # full-ish-scale A4

def quantize(x, bits):
    levels = 2**bits
    step   = 2.0 / levels         # x in [-1, 1]
    return step * np.round(x / step)

for b in (16, 8, 4, 2):
    xq  = quantize(x, b)
    err = xq - x
    snr = 10*np.log10(np.mean(x**2) / np.mean(err**2))
    print(f"{b:2d}-bit  measured SNR = {snr:6.2f} dB")
\end{lstlisting}
You will see the measured SNR track $6.02b+1.76$ closely for the higher bit depths, and the error
audibly turns from gentle hiss (high bits) into gritty, signal-correlated distortion (2--4 bits).
That correlation -- the error ``following'' the tone -- is the failure of the random model at low
bit depth, and the reason dither exists.
\end{labbox}

\section{How a real ADC works}
\label{ch19:adc-arch}

So far the quantizer has been an abstract staircase. How does silicon actually find the nearest
rung? There are several architectures; two dominate the world you will build in, and they sit at
opposite ends of a speed/resolution trade-off.

\subsection{Successive approximation (SAR): binary search in hardware}
\label{ch19:sar}
The successive-approximation register (SAR) ADC is the standard architecture inside almost every
microcontroller. It works by binary search.

The converter holds an internal DAC and a comparator. To convert an unknown input $x$, it asks a
sequence of yes/no questions, one per bit, from the most significant bit down:
\begin{enumerate}
  \item Guess the MSB is 1: set the trial DAC output to half of full scale,
        $V_{\mathrm{ref}}/2$. Ask the comparator: is $x$ above this?
  \item If yes, keep that bit as 1; if no, clear it to 0. Either way, move to the next bit and
        add (or not) the next halving, $V_{\mathrm{ref}}/4$.
  \item Repeat down to the LSB. After $N$ comparisons the register holds the code.
\end{enumerate}
Each step halves the remaining uncertainty -- exactly binary search on the value axis. An $N$-bit
SAR takes $N$ comparator decisions, so its conversion time grows only \emph{linearly} with
resolution. That makes SAR fast, low-power, and cheap, which is why it is everywhere from 10-bit
to 18-bit and from kHz to several MHz sample rates.

\begin{figure}[H]
\centering
\begin{tikzpicture}[font=\small,>={Stealth[length=2mm]}]
\node[draw, rounded corners, minimum height=8mm, fill=blue!4] (sh)  at (0,0)   {S/H};
\node[draw, isosceles triangle, isosceles triangle apex angle=60, minimum height=10mm, fill=green!5] (cmp) at (2.8,0) {};
\node at (2.6,0) {$>$};
\node[draw, rounded corners, minimum height=8mm, fill=orange!6] (sar) at (5.6,0) {SAR logic};
\node[draw, rounded corners, minimum height=8mm, fill=violet!6] (dac) at (5.6,-1.7) {internal DAC};
\node[left=4mm of sh] (in) {$x$};
\node[right=5mm of sar] (out) {code};
\draw[->] (in) -- (sh);
\draw[->] (sh) -- (cmp);
\draw[->] (cmp) -- (sar);
\draw[->] (sar) -- (out);
\draw[->] (sar) -- (dac);
\draw[->] (dac) -| (cmp.south);
\end{tikzpicture}
\caption{A SAR ADC. A sample-and-hold (S/H) freezes the input; the SAR logic toggles bits of an
internal DAC and uses the comparator to do a binary search, settling on the code in $N$ steps.}
\label{ch19:fig-sar}
\end{figure}

\begin{whybox}
Why the sample-and-hold matters: the binary search takes $N$ comparisons during which the input
must not move, or the early (big) decisions and the late (small) decisions would be about
\emph{different} voltages and the answer would be garbage. The S/H grabs the value at one instant
and holds it steady on a capacitor while the search runs. This is the hardware that realizes the
``sample at an instant'' idealization of Chapter~18.
\end{whybox}

\begin{example}[A SAR binary search by hand]
\label{ch19:ex-sar}
A 4-bit SAR with $V_{\mathrm{ref}}=8$~V converts $x=5.3$~V. Without a half-LSB rounding offset (a plain truncating search), the bit weights are $4,2,1,0.5$~V:
\begin{itemize}
  \item Bit 3: trial $4.0$. Is $5.3>4.0$? Yes $\to$ bit $=1$, accumulator $=4.0$.
  \item Bit 2: trial $4.0+2.0=6.0$. Is $5.3>6.0$? No $\to$ bit $=0$, accumulator stays $4.0$.
  \item Bit 1: trial $4.0+1.0=5.0$. Is $5.3>5.0$? Yes $\to$ bit $=1$, accumulator $=5.0$.
  \item Bit 0: trial $5.0+0.5=5.5$. Is $5.3>5.5$? No $\to$ bit $=0$, accumulator stays $5.0$.
\end{itemize}
Code $=1010_2 = 10$, representing $5.0$~V. The true value was $5.3$~V; the residual error is
$-0.3$~V. Note $\Delta = 0.5$~V here, so $\Delta/2 = 0.25$~V; the $0.3$~V residual exceeds it
because the simple truncating search lands on the rung \emph{below}. A real SAR adds the half-LSB
offset to its comparator threshold so the worst-case error becomes the promised $\Delta/2$.
\end{example}

\subsection{Sigma-delta: trade speed for resolution by oversampling}
\label{ch19:sigmadelta}
The other architecture you will meet -- in audio codecs, in 24-bit sensor front-ends, in the ADCs
on precision IMU and EMG chips -- is the \textbf{sigma-delta} ($\Sigma\Delta$) converter. It works
almost the opposite way from SAR. Where SAR resolves a value finely in one shot, sigma-
delta uses a crude quantizer (often just \emph{one bit}) running very fast, then does the real
work in arithmetic.

The method has three ingredients.

\paragraph{1. Oversampling spreads the noise out.} Suppose you sample far faster than Nyquist
requires -- by an \textbf{oversampling ratio} $\mathrm{OSR} = f_s / (2 f_{\max})$. The total
quantization noise power is still $\Delta^2/12$, but now it is spread (white) across a much wider
band, from $0$ to $f_s/2$. The fraction that falls in your actual signal band $[0, f_{\max}]$ is
only $1/\mathrm{OSR}$ of the total. After you digitally low-pass filter and discard everything
above $f_{\max}$, the in-band noise power drops by $\mathrm{OSR}$:
\[
  \sigma_{\mathrm{in\text{-}band}}^2 \;=\; \frac{\Delta^2/12}{\mathrm{OSR}}.
\]
In decibels that is a gain of $10\log_{10}\mathrm{OSR}$. Every quadrupling of the sample rate
($\mathrm{OSR}\times 4$) wins $\approx 6$~dB -- one bit of effective resolution -- \emph{for free},
purely by sampling faster and filtering.

\begin{figure}[H]
\centering
\begin{tikzpicture}
\begin{axis}[
  width=12cm, height=5cm,
  title={Oversampling spreads noise across wider bandwidth},
  title style={font=\small},
  xlabel={frequency}, ylabel={noise power density},
  xmin=0, xmax=1, ymin=0, ymax=0.5,
  axis lines=left, enlargelimits=false,
  xtick={0,0.1,1}, xticklabels={$0$,$f_{\max}$,$f_s/2$},
  legend pos=north east, legend cell align=left,
]
\addplot[draw=none, fill=blue!12, domain=0:1, samples=2] {0.25} \closedcycle;
\addlegendentry{noise spread to $f_s/2$}
\addplot[draw=none, fill=specred!35, domain=0:0.1, samples=2] {0.25} \closedcycle;
\addlegendentry{in-band ($<f_{\max}$)}
\addplot[specred, thick, domain=0:1, samples=2] {0.25};
\end{axis}
\end{tikzpicture}
\caption{Oversampling. The same total quantization noise (area under the floor) is spread across a
much wider band, so the slice that lands in the signal band $[0,f_{\max}]$ (red) is small. A
digital filter keeps the red part and discards the rest.}
\label{ch19:fig-oversample}
\end{figure}

\paragraph{2. Noise shaping pushes the noise uphill.} Sigma-delta does much better than plain
oversampling by \emph{shaping} the noise spectrum so it slopes steeply upward with frequency,
piling almost all of the quantization noise into the high band that the filter will throw away. Write one discrete-time loop step: an accumulator integrates the difference (the ``delta'') between input $x[n]$ and previous 1-bit output $y[n-1]$, giving $w[n] = w[n-1] + x[n] - y[n-1]$, and the coarse quantizer adds error $e[n]$ so $y[n] = w[n] + e[n]$. Substitute $w[n-1] = y[n-1] - e[n-1]$ and $y[n-1]$ cancels out:
\[
  y[n] \;=\; x[n] + \bigl(e[n] - e[n-1]\bigr), \qquad Y(z) \;=\; X(z) + (1 - z^{-1})\,E(z).
\]
The signal $X(z)$ passes cleanly with gain $1$, while the quantization noise $E(z)$ is multiplied by the first-difference high-pass $1 - e^{-\jj\omega}$ (magnitude $2\abs{\sin(\omega/2)} \approx \omega$ near DC), starving low frequencies of noise and piling it up near $f_s/2$. A first-order loop wins $9$~dB per octave of OSR instead of $3$; higher-order loops win even more. This is
why a 1-bit modulator at a few MHz can deliver 20+ effective bits in the audio band.

\begin{figure}[H]
\centering
\begin{tikzpicture}[font=\small,>={Stealth[length=2mm]}]
\node (in) {$x$};
\node[draw, circle, right=8mm of in] (sum) {$+$};
\node[draw, rounded corners, right=10mm of sum, fill=blue!4] (int) {$\int$};
\node[draw, rounded corners, right=10mm of int, fill=green!5] (q) {1-bit Q};
\node[right=10mm of q] (out) {bitstream};
\node[draw, rounded corners, below=10mm of int, fill=violet!6] (dac) {1-bit DAC};
\draw[->] (in) -- (sum);
\draw[->] (sum) -- (int);
\draw[->] (int) -- (q);
\draw[->] (q) -- (out);
\draw[->] (q.south) |- (dac.east);
\draw[->] (dac.west) -| node[near end, left, font=\footnotesize] {$-$} (sum.south);
\end{tikzpicture}
\caption{A first-order sigma-delta modulator. The integrator accumulates the error between input
and the fed-back 1-bit estimate; the feedback drives the running average of the output to equal
the input, while pushing quantization noise to high frequency where a decimation filter removes
it.}
\label{ch19:fig-sigmadelta}
\end{figure}

\paragraph{3. Decimation turns the fast bitstream into clean samples.} The modulator emits a
stream of single bits at the high oversampled rate. A digital low-pass \emph{decimation} filter
averages and downsamples this stream to the final output rate, producing high-resolution
multi-bit samples in the signal band. (Decimation gets its own treatment in Chapter~22.) The net
effect: the analog part is trivial (a comparator and an integrator), and most of the work is
done by cheap digital arithmetic, which suits modern chips.

\begin{focusbox}
\textbf{SAR vs.\ sigma-delta at a glance.} SAR: one sample, $N$ fast comparisons, moderate
resolution, easy to multiplex across many channels, low latency -- ideal for control loops, IMUs,
general MCU sensing. Sigma-delta: oversample hard, shape the noise, decimate -- very high
resolution, built-in anti-alias filtering, but higher latency and best for one continuous
channel -- ideal for audio and precision instrumentation.
\end{focusbox}

\section{Dithering: adding noise to get resolution}
\label{ch19:dither}

Return to the failure mode we keep flagging: when the signal is small or slowly varying, the
quantization error stops looking random. It becomes \emph{correlated} with the signal, and
correlated quantization error produces audible \textbf{harmonic distortion} rather than benign white hiss: harmonics that the ear and the
spectrum analyzer both find ugly. The cure is counterintuitive: \textbf{add a small
amount of random noise before quantizing.}

\begin{intuitionbox}
Imagine a quiet sine wave whose entire swing fits between two rungs of the ladder. Without help,
the quantizer reports a flat line at the nearer rung -- the tone vanishes into a stair-step and
its information is lost. Now sprinkle in a little noise of about one LSB. On each sample the noisy
value sometimes pushes up to the next rung and sometimes down. The \emph{individual} outputs are
wrong, but their \emph{running average} tracks the true sub-LSB value. You have traded a coherent,
audible error for incoherent hiss -- and incoherent hiss is far easier to ignore and can be
averaged or filtered away.
\end{intuitionbox}

\begin{whybox}
Dither works because averaging recovers information that a single quantized sample cannot hold.
A value sitting $0.3$~LSB above a rung, dithered, will round up about $30\%$ of the time and down
$70\%$ of the time; the long-run mean of the codes is exactly $0.3$~LSB above the rung. The dither
\emph{linearizes} the quantizer on average: it converts a deterministic, signal-dependent rounding
into an unbiased estimator whose error is independent of the signal. You pay a small, fixed
increase in noise floor and in return you remove the distortion harmonics. For
well-chosen dither --- \textbf{TPDF dither} (triangular PDF spanning $[-\Delta,+\Delta]$, formed simply by summing two independent uniform $[-\Delta/2,+\Delta/2]$ random numbers) --- the residual error is not only
decorrelated but its variance becomes independent of the signal level.
\end{whybox}

Dithering is standard well beyond audio: in high-quality ADCs and DACs, in image
processing (the ``dithered'' gradients that hide banding in 8-bit images), and in any place a coarse
grid would otherwise create visible or audible structure. The general idea -- \emph{break correlated
error into uncorrelated noise, then average} -- recurs throughout signal processing.

\begin{applicationbox}
\textbf{Audio mastering.} When a 24-bit studio mix is reduced to 16-bit for a CD, engineers add
TPDF (triangular probability density) dither, often with a spectral shape that pushes the added
noise into frequencies where the ear is least sensitive (``noise-shaped dither''). The result is
that quiet reverb tails fade smoothly into silence instead of chopping off in audible steps, even
though only 16 bits are stored.
\end{applicationbox}

\section{From bits back to the world: the DAC}
\label{ch19:dac}

A digital-to-analog converter runs the staircase in reverse: given a code, produce a voltage. The
simplest mental model is a bank of switches and a resistor network (an ``R-2R ladder'') that adds up
weighted contributions of each bit. But the geometry that matters is what comes \emph{out}.

A DAC does not emit the ideal samples (infinitely thin pins). It holds each sample's voltage
constant until the next one arrives -- a \textbf{zero-order hold} (ZOH). The output is a
staircase in time, not a smooth curve. That staircase is the samples convolved with a rectangular
pulse one sample wide, and convolution with a rectangle in time means multiplication by a
\emph{sinc} in frequency:
\[
  H_{\mathrm{ZOH}}(f) \;=\; T_s\,\sinc(f T_s)\, e^{-\jj \pi f T_s},
  \qquad \abs{H_{\mathrm{ZOH}}(f)} = T_s\,\abs{\sinc(f T_s)},
\]
where the width-$T_s$ pulse gives the $T_s\sinc(f T_s)$ magnitude and its midpoint at $t=T_s/2$ gives the half-sample linear phase lag $e^{-\jj 2\pi f(T_s/2)} = e^{-\jj\pi f T_s}$. Two consequences follow, and both have fixes.

\begin{enumerate}
  \item \textbf{High-frequency droop.} The sinc magnitude sags as $f$ rises, reaching about
        $-3.9$~dB at the band edge $f = f_s/2$. Your highest wanted frequencies come out slightly
        quiet. Fix: a gentle digital or analog \emph{de-emphasis} that boosts the top of the band
        by the inverse sinc.
  \item \textbf{Spectral images.} The staircase contains copies (images) of your spectrum around
        every multiple of $f_s$. Fix: an analog \textbf{reconstruction filter} (low-pass) after
        the DAC that smooths the steps and removes the images -- the mirror image of the
        anti-alias filter before the ADC.
\end{enumerate}

\begin{figure}[H]
\centering
\begin{tikzpicture}
\begin{axis}[
  width=12cm, height=5cm,
  title={DAC zero-order hold output waveform},
  title style={font=\small},
  xlabel={time}, ylabel={DAC output},
  xmin=0, xmax=1, ymin=-1.25, ymax=1.25,
  axis lines=left, enlargelimits=false,
  legend pos=south east, legend cell align=left,
]
\addplot[sigblue, thick, domain=0:1, samples=200] {sin(deg(2*pi*1.5*x))};
\addlegendentry{ideal signal}
\addplot[specred, thick, const plot, no marks] coordinates {
  (0,0)(0.1,0.81)(0.2,0.95)(0.3,0.31)(0.4,-0.59)(0.5,-1.0)
  (0.6,-0.59)(0.7,0.31)(0.8,0.95)(0.9,0.81)(1.0,0.0)
};
\addlegendentry{zero-order hold}
\end{axis}
\end{tikzpicture}
\caption{A DAC's zero-order-hold output (red) holds each sample until the next. The steps contain
high-frequency images and a slight sinc droop; an analog reconstruction filter smooths them back
toward the ideal curve (blue).}
\label{ch19:fig-zoh}
\end{figure}

\begin{connectionbox}
The ZOH droop comes from the same $\rect \leftrightarrow \sinc$ duality we saw in Chapter~18. There, sampling in time put
a sinc-shaped interpolation kernel between samples; here, holding in time multiplies the spectrum
by a sinc. Same rectangle, same transform pair $\rect \leftrightarrow \sinc$ -- just applied on
the way out instead of the way in. Whenever you see a sinc, look for a rectangle nearby.
\end{connectionbox}

\section{Reading a real sensor through an ADC}
\label{ch19:bench}

On real hardware, the theory above reduces to a handful of practical numbers.

\begin{buildbox}
\textbf{Arduino Uno (ATmega328P): 10-bit SAR ADC.} The Uno has one SAR converter multiplexed
across six pins. Key facts:
\begin{itemize}
  \item Resolution $N=10 \Rightarrow 1024$ codes. With the default $V_{\mathrm{ref}}=5$~V,
        $\Delta = 5/1024 \approx 4.9$~mV per code. You cannot resolve a change smaller than that.
  \item \texttt{analogRead()} returns $0\ldots1023$. Convert to volts with
        \texttt{v = raw * 5.0 / 1024.0}.
  \item Default conversion time is about $100\,\mu s$ ($\sim$9600 reads/s), set by the ADC clock
        prescaler. Plenty for an IMU at 100--1000~Hz; far too slow for audio.
  \item For quiet measurements use the internal 1.1~V reference (\texttt{analogReference(INTERNAL)}):
        now $\Delta \approx 1.07$~mV, about four times finer, at the cost of a smaller input range.
\end{itemize}
\begin{lstlisting}[language=C]
// Read a flex sensor / EMG envelope on A0, oversample to fake extra bits.
void setup() { Serial.begin(115200); }

void loop() {
  long acc = 0;
  const int OS = 16;             // average 16 reads
  for (int i = 0; i < OS; i++) acc += analogRead(A0);
  // OSR=16 -> +2 effective bits IF the signal jitters by >= 1 LSB
  float code12 = acc / 4.0;      // 10-bit*16 -> ~12-bit scale
  float volts  = code12 * (5.0 / 4096.0);
  Serial.println(volts, 4);
  delay(5);                      // ~200 Hz output
}
\end{lstlisting}
The averaging only buys real resolution if the input (plus its natural noise) wiggles by at least
an LSB between reads -- that is dithering happening for free. On a dead-quiet, perfectly stable
voltage, averaging 16 identical reads gives you back the same number and no extra bits.
\end{buildbox}

\begin{buildbox}
\textbf{Raspberry Pi: no analog pins at all.} The Pi's GPIO is digital only. To read an analog
sensor you add an external ADC over SPI or I\textsuperscript{2}C:
\begin{itemize}
  \item \textbf{MCP3008}: 8-channel, 10-bit SAR, SPI, up to $\sim$200~kS/s. Cheap and ubiquitous;
        the standard ``give my Pi an ADC'' chip.
  \item \textbf{ADS1115}: 4-channel, 16-bit sigma-delta, I\textsuperscript{2}C, with a
        programmable-gain amplifier. Slower (up to 860~S/s) but far finer -- a good match for an
        EMG envelope or any low-frequency, high-dynamic-range signal.
\end{itemize}
\begin{lstlisting}[language=Python]
# Raspberry Pi + ADS1115 (16-bit sigma-delta) over I2C
import board, busio
import adafruit_ads1x15.ads1115 as ADS
from adafruit_ads1x15.analog_in import AnalogIn

i2c = busio.I2C(board.SCL, board.SDA)
ads = ADS.ADS1115(i2c)
ads.gain = 1                      # +/-4.096 V full scale
chan = AnalogIn(ads, ADS.P0)      # single-ended on A0

while True:
    print(f"{chan.voltage:7.4f} V  (raw {chan.value})")
\end{lstlisting}
With $\pm4.096$~V full scale and 16 bits, $\Delta \approx 8.192/65536 \approx 125\,\mu$V -- about
40 times finer than the Uno's default. The sigma-delta front-end also gives you built-in
anti-alias filtering, which matters for the low, clean signals (EMG/FMG envelopes) we want.
\end{buildbox}

Figure~\ref{ch19:fig-emg} shows why the converter choice matters here. An EMG envelope is a
small, slow bump riding on a steady baseline; if the ADC's $\Delta$ is comparable to the bump's
height, the staircase swallows the feature whole. A finer step traces the same bump smoothly.

\begin{figure}[H]
\centering
\begin{tikzpicture}
\begin{axis}[
  width=12cm, height=5cm,
  title={Coarse vs.\ fine quantization on a bio-signal feature},
  title style={font=\small},
  xlabel={time (s)}, ylabel={EMG envelope (V)},
  xmin=0, xmax=2, ymin=0, ymax=2,
  axis lines=left, domain=0:2, samples=200,
  enlargelimits=false, legend pos=north east, legend cell align=left,
]
% slow baseline plus a gentle activation bump
\addplot[sigblue, thick, samples=200] {0.6 + 0.7*exp(-((x-1.0)^2)/0.06)};
\addlegendentry{true envelope}
% coarse quantization: step 0.25 V (close to bump height) loses detail
\addplot[specred, thick, const plot, no marks, samples=40]
  {0.25*round((0.6 + 0.7*exp(-((x-1.0)^2)/0.06))/0.25)};
\addlegendentry{coarse $\Delta=0.25$~V}
\end{axis}
\end{tikzpicture}
\caption{A slow EMG-style envelope (blue) and the same signal read by a converter whose step
$\Delta$ is a large fraction of the feature height (red staircase). Coarse quantization smears the
activation into a few plateaus; a 16-bit front-end with a small $\Delta$ keeps the shape.}
\label{ch19:fig-emg}
\end{figure}

\begin{goalbox}
For the wrist controller, match the converter to the signal. The \textbf{IMU} (accelerometer +
gyro, Chapter~34) already contains its own high-quality sigma-delta ADCs and hands you clean
multi-bit samples over I\textsuperscript{2}C/SPI -- you never touch its quantizer, but its ENOB
and noise density (specified in $\mu g/\sqrt{\mathrm{Hz}}$) set the floor on how small a wrist
tilt you can detect. The \textbf{muscle signals} that hint at finger intent (EMG/FMG,
Chapter~37) are tiny, low-frequency, and ride on a large slow baseline -- a textbook case for a
16-bit sigma-delta front-end like the ADS1115 with a PGA, where dithering and oversampling
genuinely buy you resolution. And keep a clear eye on what bits can and cannot do: a finer ADC reads
the forearm's coarse muscle activity more precisely, but \emph{no} bit depth lets a wristband
recover continuous individual finger-joint angles -- that information is not present in the signal
the wristband sees (Chapter~37 makes this limit precise). Finer quantization captures the signal more
cleanly but cannot add information the signal lacks.
\end{goalbox}

\begin{historybox}
When digital communications and pulse-code modulation (PCM) took off in the mid-20th century,
engineers faced a basic constraint: continuous physical voltages cannot be stored
in finite digital memory without rounding. Alec Reeves invented PCM in 1937, patenting the idea of
encoding speech into binary pulses, but the vacuum tubes of the 1930s were far too bulky and power-hungry
to make commercial converters practical. Only with the advent of transistors and monolithic silicon
in the 1960s did fast, accurate analog-to-digital conversion become routine.

A persistent puzzle was how to model the resulting rounding error. Is quantization a severe,
signal-dependent nonlinear distortion, or can it be treated as benign random noise? Bernard Widrow
at Stanford put this on rigorous footing in 1956 with his Quantization Theorem, proving that if an
input signal fluctuates across multiple quantization steps, the error behaves almost exactly like
independent uniform white noise of variance $\Delta^2/12$. Meanwhile, dither had an accidental wartime
origin: World War II aviators discovered that mechanical analog bomb-sights calculated target trajectories
more accurately in vibrating aircraft than resting on a test bench. The airframe vibrations were
unintentionally dithering the gears, nudging them past mechanical friction. Postwar engineers deliberately
injected electrical noise into quantizers to eliminate harmonic distortion.
\end{historybox}

\section{Summary}
\label{ch19:summary}

Quantization is the vertical twin of sampling. An $N$-bit ADC snaps each sample to the nearest of
$2^N$ rungs spaced $\Delta = V_{\mathrm{ref}}/2^N$ apart; the error is bounded by $\pm\Delta/2$
and, for busy signals, behaves like uniform noise of power $\Delta^2/12$. Comparing that fixed
noise to a full-scale sinusoid gives the rule
$\mathrm{SNR}\approx 6.02N + 1.76$~dB -- about a bit per 6~dB -- and inverting it on a real
converter gives its ENOB. SAR converters do a fast hardware binary search for moderate-resolution,
low-latency, multi-channel sensing; sigma-delta converters oversample a crude quantizer, shape the
noise out of band, and decimate, reaching very high resolution for one continuous channel.
Oversampling alone buys $\approx 3$~dB per doubling of $f_s$; noise shaping does much better; and
\emph{adding} noise (dither) converts ugly signal-correlated distortion into benign hiss you can
average away. On the way back out, a DAC holds each sample (zero-order hold), which droops the top
of the band and creates spectral images that a reconstruction filter cleans up. At the bench, the
Arduino gives you a 10-bit SAR for free; the Pi needs an external ADC, where a 16-bit sigma-delta
part is the right tool for small, slow biosignals -- always remembering that resolution captures
what the signal contains and never more.

Having quantified the noise created by discretization itself, we now turn to the noise already
present in the physical world. Chapter~20 investigates thermal agitation, random processes,
autocorrelation, and how to pull weak signals from noisy backgrounds.

\section{Exercises}
\label{ch19:exercises}

\begin{exercise}
An ADC has $V_{\mathrm{ref}}=3.3$~V and $N=12$ bits. (a) Compute the LSB size $\Delta$ in
millivolts. (b) What is the maximum quantization error in millivolts? (c) What is the ideal SNR in
dB for a full-scale sinusoid?
\end{exercise}

\begin{exercise}
Derive the quantization noise power $\sigma_e^2 = \Delta^2/12$ from scratch by integrating $e^2$
against a uniform density on $[-\Delta/2,\Delta/2]$. Then state in one sentence why the answer is
independent of where the signal sits on the ladder.
\end{exercise}

\begin{exercise}
You record audio at 16 bits but, through a gain mistake, your loudest peaks only reach
$1/8$ of full scale. (a) By how many dB has your signal power dropped relative to full scale?
(b) Estimate your effective number of bits. (c) State the practical lesson.
\end{exercise}

\begin{exercise}
A SAR ADC has $N=8$ bits and $V_{\mathrm{ref}}=4.096$~V (so $\Delta = 16$~mV). Trace the binary
search by hand for an input of $2.7$~V: list each trial voltage, each comparator decision, and the
final 8-bit code. What voltage does that code represent, and what is the residual error?
\end{exercise}

\begin{exercise}
You oversample a signal by $\mathrm{OSR}=64$ (no noise shaping) and then ideally low-pass filter
to the signal band. (a) How many dB of in-band SNR improvement do you gain? (b) How many effective
bits is that? (c) Why does this require the quantization error to be uncorrelated across samples,
and what tool guarantees that?
\end{exercise}

\begin{exercise}
Explain, in your own words and with a small numeric example, how dithering lets the long-run
\emph{average} of quantized codes represent a value that lies strictly between two rungs. Use a
true value of $0.4$~LSB above a rung and predict the fraction of samples that round up.
\end{exercise}

\begin{exercise}
A DAC runs at $f_s = 48$~kHz with a zero-order hold. (a) Write the magnitude of the ZOH frequency
response. (b) Compute the droop in dB at $f = 20$~kHz relative to DC. (c) Describe the two analog
problems the ZOH creates and the filter that fixes each.
\end{exercise}

\begin{exercise}
Modify the Arduino sketch in Section~\ref{ch19:bench} to switch to the internal 1.1~V reference.
(a) What is the new $\Delta$? (b) Under what condition on your sensor's output range is this a good
idea, and when would it clip? (c) Why does oversampling-and-averaging fail to add bits on a
perfectly stable, noise-free input, and what is the one-sentence connection to dither?
\end{exercise}

\begin{exercise}
A ``16-bit'' ADC datasheet quotes a measured SNR of 91~dB. (a) Compute its ENOB. (b) The same part
quoted with a $4\times$ slower (more averaged) output mode reports 97~dB; compute the new ENOB and
explain physically where the extra bits came from. (c) Why can no amount of averaging push the
ENOB past the analog noise floor of the front-end?
\end{exercise}

\begin{exercise}
For the wrist controller, you must choose an ADC for an EMG envelope signal that spans roughly
0--2~V, changes slowly (under 50~Hz of useful bandwidth), and must resolve features as small as a
few hundred microvolts. (a) Would you pick a 10-bit SAR at 200~kS/s or a 16-bit sigma-delta at
860~S/s? Justify using $\Delta$ and the bandwidth. (b) State one thing this ADC choice still
\emph{cannot} tell you about the user's individual fingers, and why.
\end{exercise}
